Title: Formalising Razborov's Flag Algebras in Lean

Abstract: 
Razborov's flag algebras are a powerful method for proving asymptotic inequalities in
extremal combinatorics. They often reduce an asymptotic extremal problem over combinatorial objects to a
semidefinite-programming (SDP) problem, and in doing so, they enable the automatic resolution 
of the problem by a computer. In this talk, I will explain our experience of formalising flag algebras 
for finite simple graphs in Lean for the past year and a half. My colleagues and I very much enjoyed 
writing Lean code for flag algebras, but at the same time, we were constantly worried that others 
might not care about the formalisation as much as we hoped, because flag algebras are believed to
have no major errors. In the talk, I will also describe our attempts to address this concern. Using the meta-programming
capabilities of Lean and the reflection mechanism, we built a semi-automatic pipeline for converting 
an SDP certificate of a flag-algebra proof to a formal proof in Lean. We also developed a meta-theory that
justifies the deviation of our formalisation from the original flag algebras, which was needed to make our 
formalisation project more tractable and manageable. This meta-theory itself later was formalised in Lean as well,
but automatically by Claude due to the formalisation that we built manually. 

This is joint work with Jihoon Hyun, Gyeongwon Jeong, Sang-il Oum, and Seonghun Park.
